\subsection{Specialized Workflow vs.\ General-Purpose Agent (RQ2)}
\label{subsec:general_agent}

We next examine whether SoCPilot's structured hardware/SoC workflow is more reliable than open-ended execution by a general-purpose LLM agent. Following HSCO-Bench, the general agent is given the application and the available SoC development environment and attempts to progress through accelerator generation and the downstream tool flow. The comparison focuses on the shared automated stages and does not count manually authored embedded runtime software as a SoCPilot capability.

\begin{table*}[t]
\centering
\caption{Execution results of the general-purpose agent following the HSCO-Bench protocol on our application set.}
\label{tab:hscobench_results}
\renewcommand{\arraystretch}{1.18}
\setlength{\tabcolsep}{4.5pt}
\resizebox{\textwidth}{!}{
\begin{tabular}{l l r l l p{6.0cm}}
\toprule
\textbf{Application} &
\textbf{Generated Accelerator Scope} &
\textbf{\begin{tabular}[c]{@{}c@{}}Agent Time\\(s)\end{tabular}} &
\textbf{Furthest Passed Stage} &
\textbf{Failure Stage} &
\textbf{Failure Cause} \\
\midrule
\texttt{nightvision} &
Laplacian fusion (\texttt{lapfusion}) &
3974 &
Device-ID check &
BEV-Sim compilation &
Invalid C++ type operation: \texttt{to\_uint()} is applied to variables of type \texttt{uint32\_t}. \\

\texttt{scan} &
Scan post-processing (\texttt{scanpp}) &
2563 &
BEV-Sim compilation/execution &
BEV-Sim verification &
Functional mismatch with the software golden output; 17 output pixels differ from the reference results. \\

\texttt{plate\_recognition} &
Plate-localization pipeline (\texttt{plateaccel}) &
2325 &
Device-ID check &
BEV-Sim compilation &
Invalid C++ type operation: \texttt{to\_int()} is applied to configuration fields of type \texttt{int32\_t}. \\

\texttt{recognition} &
HSV mask and dilation pipeline (\texttt{recog}) &
1910 &
BEV-Sim &
Catapult HLS &
Variable-division operators cannot be scheduled under the available Catapult resource and multicycle constraints. \\

\texttt{wami} &
WAMI frame pipeline (\texttt{wami}) &
4024 &
BEV-Sim &
Catapult HLS &
Catapult reports no definition for \texttt{floorf()} during HLS compilation. \\
\bottomrule
\end{tabular}}
\end{table*}

The general-agent failures span type correctness, functional equivalence, and HLS compatibility. To make this comparison visually explicit, the final paper should include a stage-pass matrix showing, for each application and method, whether it passes accelerator-code generation, functional simulation, HLS synthesis, and ESP integration. FPGA board execution should be shown for SoCPilot as an RQ1 validation result, not as a common RQ2 comparison stage unless the baseline is also evaluated through the same board stage.

% Recommended figure (data/status still required):
% Fig. Y. Stage-pass matrix for General Agent vs. SoCPilot.
% Columns: Accelerator generation | Functional simulation | HLS | ESP integration.
% Optional SoCPilot-only annotation: FPGA board validated (reported under RQ1).

